% Part: first-order-logic
% Chapter: model-theory
% Section: nonstandard-arithmetic

\documentclass[../../../include/open-logic-section]{subfiles}

\begin{document}

\olfileid{mod}{bas}{nsa}
\section{అంకగణితపు అప్రమాణ నమూనాలు}

\begin{defn}
$\Lang{L_N}$ అనేది అంకగణితపు \tetoken{భాష}{!!{language}} అనుకుందాం. దీనిలో \tetoken{స్థిరాంకం}{!!{constant}} $\Obj{0}$, ద్విస్థానిక \tetoken{విధేయం}{!!{predicate}} $<$, ఏకస్థానిక \tetoken{ప్రమేయ సంకేతం}{!!{function}} $\prime$, ద్విస్థానిక \tetoken{ప్రమేయ సంకేతాలు}{!!{function}s} $+$, $\times$ ఉంటాయి.
\begin{enumerate}
\item అంకగణితపు \emph{ప్రమాణ నమూనా} అంటే $\Lang{L_N}$ కోసం \tetoken{నిర్మాణం}{!!{structure}} $\Struct{N}$. దీనిలో $\Nat = \{ 0, 1, 2, \dots\}$; $\Obj{0}$ను $0$గా, $<$ను~$\Nat$పై చిన్నదనే సంబంధంగా, $\prime$, $+$, $\times$లను వరుసగా $\Nat$పై ఉత్తరవర్తి, కూడిక, గుణకారాలుగా అన్వయిస్తాం.
\item \emph{సత్య అంకగణితం} అంటే సిద్ధాంతం $\Theory{N}$.
\end{enumerate}
\end{defn}

$\Lang{L_N}$లో పనిచేసేటప్పుడు, $\Obj{0}$కు ఉత్తరవర్తి ప్రమేయాన్ని $n$సార్లు వర్తింపజేసిన $\Obj{0}^{\prime\dots\prime}$ రూపంలోని ప్రతి పదాన్ని~$\num{n}$ అని సంక్షిప్తంగా రాస్తాం.

\begin{defn}
$\Lang{L_N}$ కోసం \tetoken{నిర్మాణం}{!!{structure}}~$\Struct{M}$ \emph{ప్రమాణమైనది} కావడానికి, కావలసినదీ చాలునదీ $\Struct{N} \iso \Struct{M}$.
\end{defn}

\begin{thm}\ollabel{thm:non-std}
సత్య అంకగణితానికి \tetoken{గణనీయ}{!!{enumerable}} అప్రమాణ నమూనాలు ఉన్నాయి.
\end{thm}

\begin{proof}
$\Lang{L_N}$లో కొత్త \tetoken{స్థిరాంకం}{!!{constant}}~$c$ను చేర్చి, ఈ సిద్ధాంతాన్ని పరిశీలించండి:
\[
\Theory{N} \cup \Setabs{\num{n} < c}{n \in \Nat}.
\]
ఈ సిద్ధాంతం పరిమిత సంతృప్త్యర్హమైనది. కనుక సంహతత్వం ప్రకారం దీనికి $\Struct{M}$ అనే నమూనా ఉంది; దిగువ లోవెన్‌హైమ్--స్కోలెమ్ సిద్ధాంతం ప్రకారం దాన్ని \tetoken{గణనీయమైనదిగా}{!!{enumerable}} తీసుకోవచ్చు. $\Domain{M}$ అనేది $\Struct{M}$ ప్రదేశం అయితే, $\Lang{L}$లోని అతార్కిక స్థిరాంకాలను $\Struct{M}$ ఈ విధంగా అన్వయిస్తుందనుకుందాం: $\mathbf{z} = \Assign{\Obj{0}}{M} \in \Domain{M}$, ${\prec} =
\Assign{<}{M} \subseteq M^2$, $* = \Assign{\prime}{M} \colon
\Domain{M} \to \Domain{M}$; అలాగే $\oplus = \Assign{+}{M}, \otimes =
\Assign{\times}{M} \colon \Domain{M}^2 \to \Domain{M}$. ప్రతి $x
\in \Domain{M}$కు $*$ను వర్తింపజేసి లభించే $\Domain{M}$ మూలకాన్ని $x^*$గా రాస్తాం.
\sourcecorrection{OLTEMODBASENSA-001}{ఈ నిరూపణ మొదట $\Lang{L_N}$ను విస్తరించి, వెంటనే దాని మూల సంకేతాల అన్వయాన్ని $\Lang{L}$ది అని పిలుస్తుంది. అదే అంకగణిత భాష ఉద్దేశమని సందర్భం సూచించినా, మూల సంకేత అసంగతిని ప్రకటించాం.}

ఇప్పుడు $\Struct{N}$, $\Struct{M}$ల మధ్య $h$ ఒక సమరూపత అయితే, $h(n) = \Assign{c}{M}$ అయ్యే $n \in \Nat$ ఒకటి ఉంటుంది. కాబట్టి $\Struct{N}$లో $s(x) = n$ అయ్యే ఏ కేటాయింపు $s$నైనా తీసుకోండి. అప్పుడు $\Sat{N}{\eq[\num{n}][x]}[s]$; \olref[iso]{thm:isom} నిరూపణ ప్రకారం $\Sat{M}{\eq[\num{n}][x]}[h \circ s]$ కూడా. అందువల్ల $\Assign{c}{M} =
\mathbf{z}^{*\cdots *}$ ($*$ను $n$సార్లు వర్తింపజేసినప్పుడు). కానీ పరికల్పన ప్రకారం $\Sat{M}{\num{n} < c}$, $\prec$ అస్వావర్తనమైనది కనుక ఇది అసాధ్యం. కాబట్టి $\Struct{M}$ అప్రమాణ నమూనా.
\end{proof}

\begin{prob}
సమితి $X$పై సంబంధం $R$ \emph{సుసంస్థాపితమైనది} కావడానికి, కావలసినదీ చాలునదీ~$R$లో అనంత అవరోహణ శృంఖలాలు లేకపోవడం. అంటే $X$లో $\dots x_2Rx_1Rx_0$ అయ్యే $x_0$, $x_1$, $x_2$,~\dots లేకపోవడం. జెర్మెలో--ఫ్రెంకెల్ సమితి సిద్ధాంతం~$ZF$ అవైరుధ్యమైనదని అనుకొని, $\dots x_2 \in x_1 \in x_0$ అయ్యే $\mathfrak{M}$ వంటి $ZF$ యొక్క అసుసంస్థాపిత నమూనాలు ఉన్నాయని చూపండి.
\end{prob}

\olref{thm:non-std}లోని అప్రమాణ నమూనా $\Struct{M}$, ప్రమాణ నమూనాతో ప్రాథమిక తుల్యత కలిగి ఉంది. కనుక $\Struct{M}$ గురించి కొన్ని ధర్మాలను వ్యుత్పాదించవచ్చు. ఈ విభాగం మిగతా భాగంలో ఆ పనిని చేస్తాం; దాంతో $\Theory{N}$ యొక్క \tetoken{గణనీయ}{!!{enumerable}} అప్రమాణ నమూనాలను కచ్చితంగా వర్ణించగలం.

\begin{enumerate}
\item $\Domain{M}$లో ఏ మూలకమూ తనకన్నా తాను $\prec$-చిన్నది కాదు: $\lforall[x][\lnot x < x]$ వాక్యం $\Struct{N}$లో సత్యం కనుక $\Struct{M}$లోనూ సత్యమే.
\item ఇదే విధంగా $\prec$ అనేది $\Domain{M}$పై \emph{రేఖీయ క్రమం}, అంటే $\Domain{M}$పై సంపూర్ణమైన, అస్వావర్తనమైన, సంక్రామకమైన సంబంధం అని తెలుస్తుంది.
\item $\mathbf{z}$ అనేది $\Domain{M}$లో $\prec$-అత్యల్ప మూలకం.
\item $\Domain{M}$లోని ప్రతి మూలకం తన $*$-ఉత్తరవర్తికన్నా $\prec$-చిన్నది; $x^*$ అనేది $x$కన్నా పెద్ద $\Domain{M}$ మూలకాలలో $\prec$-అత్యల్పం.
\item $\Struct{M}$లో $\Nat$కు సమరూపమైన $\prec$-ప్రారంభ ఖండం ఉంది: $\mathbf{z}, \mathbf{z}^*, \mathbf{z}^{**}, \dots$. దాన్ని $\Domain{M}$లోని \emph{ప్రమాణ భాగం} అంటాం. $\Domain{M}$లోని మిగతా ప్రతి మూలకం \emph{అప్రమాణమైనది}. $\Domain{M}$లో అలాంటి అప్రమాణ మూలకాలు తప్పక ఉండాలి; లేదంటే \olref{thm:non-std} నిరూపణలోని $h$ ప్రమేయం సమరూపత అవుతుంది. $\Struct{M}$ ప్రమాణ భాగంపై విస్తరించే \tetoken{చరరాశులుగా}{!!{variable}s} $n, m, \dots$లను వాడతాం.
\item ప్రతి అప్రమాణ మూలకం ఏ ప్రమాణ మూలకంకన్నా పెద్దది. ఎందుకంటే ప్రతి $n \in \Nat$కు
  \[
  \Sat{N}{\lforall[z][(\lnot(\eq[z][\Obj{0}] \lor
  \dots \lor \eq[z][\num{n}]) \lif \num{n} < z)]},
  \]
  కనుక $z \in \Domain{M}$ అన్ని ప్రమాణ మూలకాలకన్నా భిన్నమైతే, అది వాటన్నింటికన్నా \emph{పెద్దది}.
\item $\mathbf{z}$ కాని ప్రతి $\Domain{M}$ మూలకం, $\Domain{M}$లోని ఏదో ఒక ఏకైక మూలకానికి $*$-ఉత్తరవర్తి. దాన్ని $^*x$గా సూచిస్తాం. $x = y^*$ అయితే, వాటిలో ఒకటి ప్రమాణమైతే $x$, $y$ రెండూ ప్రమాణాలే (ఒకటి అప్రమాణమైతే రెండూ అప్రమాణాలే).
\item $\Domain{M}$పై సమానత్వ సంబంధం $\approx$ను ఇలా నిర్వచిద్దాం: ఏదో \emph{ప్రమాణ} $n$కు $x \oplus n = y$ లేదా $y \oplus n =x$ అయితే, అప్పుడు మాత్రమే $x \approx y$. అంటే $x$, $y$ల మధ్య దూరం పరిమితమైనప్పుడు, అప్పుడు మాత్రమే $x
  \approx y$. $n$, $m$ ప్రమాణాలైతే $n \approx m$. $x$ యొక్క \emph{బ్లాక్} అంటే తుల్యతా వర్గం $[x] = \Setabs{y}{x
  \approx y}$.
\item $x \prec y$ కాగా $x \not\approx y$ అనుకుందాం. $\Sat{N}{\lforall[x][\lforall[y][(x < y \lif (x' < y \lor x' =
      y))]]}$ కనుక $x^* \prec y$ లేదా $x^* = y$. రెండోది $x \approx y$ను సూచిస్తుంది కాబట్టి అసాధ్యం; అందువల్ల $x^* \prec y$. ఇదే విధంగా $x \prec y$, $x \not\approx y$ అయితే $x \prec
  {^*y}$. కాబట్టి $x \prec y$, $x \not\approx y$ అయితే ప్రతి $w \approx x$, ప్రతి $v \approx y$కన్నా $\prec$-చిన్నది. అందువల్ల ప్రతి బ్లాక్ $[x]$ రెండు వైపులా అనంతమైన ఈ శృంఖలాన్ని ఏర్పరుస్తుంది:
  \[
  \cdots \prec  {^{**}x} \prec {^*}x \prec x \prec x^* \prec x^{**}
  \prec \cdots
  \]
  దీని క్రమరకం పూర్ణ సంఖ్యల క్రమరకం కాబట్టి దీన్ని $Z$-శృంఖలమంటాం.
\sourcecorrection{OLTEMODBASENSA-002}{ఈ అంశంలో $x^*=y$ అసాధ్యమని చెప్పాక మూలం పొరపాటున మళ్లీ $x\prec y$నే తేల్చింది; అవసరమైన బలమైన ముగింపు $x^*\prec y$. తదుపరి బ్లాక్ వాదనకు దానినే స్పష్టం చేశాం.}
\item $\prec$ క్రమాన్ని బ్లాక్‌లకు ఎత్తవచ్చు: $x \prec y$ కాగా $x \not\approx y$ అయితే $x$ బ్లాక్, $y$ బ్లాక్‌కన్నా చిన్నది. బ్లాక్‌లో అప్రమాణ మూలకం ఉంటే అది \emph{అప్రమాణ} బ్లాక్. ప్రమాణ బ్లాక్ అన్నిటికన్నా చిన్నది.
\sourcecorrection{OLTEMODBASENSA-003}{మూలం ``$x\prec y$ అయితే $x$ బ్లాక్ $y$ బ్లాక్‌కన్నా చిన్నది'' అని అంది. అదే బ్లాక్‌లోనూ $x\prec y$ కావచ్చు; కాబట్టి వేర్వేరు బ్లాక్‌ల షరతు $x\not\approx y$ చేర్చాం.}
\item అత్యల్ప అప్రమాణ బ్లాక్ లేదు: $y$ అప్రమాణమైతే, $x$ కూడా అప్రమాణమై, $x
  \not\approx y$ కాగా $x \prec y$ అయ్యే $x$ ఒకటి ఉంటుంది. నిరూపణ: ప్రమాణ నమూనా $\Struct{N}$లో ప్రతి సంఖ్యను రెండు చేత భాగించవచ్చు, మిగులు ఒకటి కావచ్చు; అంటే $\Sat{N}{\lforall[y][\lexists[x][(y = x + x \lor y = x + x +
        \Obj{0}')]]}$. ప్రాథమిక తుల్యత ప్రకారం ప్రతి $y \in
  \Domain{M}$కు $x \oplus x = y$ లేదా $x \oplus x \oplus \mathbf{z}^*= y$ అయ్యే $x \in \Domain{M}$ ఒకటి ఉంటుంది. $x$ ప్రమాణమైతే $y$ కూడా ప్రమాణమవుతుంది; కనుక $x$ అప్రమాణం. అంతేకాదు, $x$, $y$ వేర్వేరు బ్లాక్‌లకు చెందుతాయి, అంటే $x \not\approx y$. దీన్ని చూడటానికి అవి ఒకే బ్లాక్‌లో ఉన్నాయనుకుందాం; అంటే ఏదో ప్రమాణ $n$కు $x \oplus n =
  y$. $y = x \oplus x$ అయితే $x \oplus n = x
  \oplus x$; కూడిక రద్దు నియమం వల్ల $x = n$ (అది $\Struct{N}$లోనూ, అందువల్ల $\Struct{M}$లోనూ వర్తిస్తుంది). అప్పుడు $x$ ప్రమాణమవుతుంది. $y = x \oplus x
  \oplus \mathbf{z}^*$ అయినప్పటికీ ఇదే తరహా వాదన.
\sourcecorrection{OLTEMODBASENSA-004}{``ప్రతి సంఖ్య రెండు చేత భాగించబడుతుంది, మిగులు ఒకటి కావచ్చు'' అనే మూల వాక్యానికి ఇచ్చిన సూత్రంలో $\lforall[x]$ ఉంది; అది అసత్యం. అవసరమైనది ప్రతి $y$కు \emph{ఏదో} $x$ ఉండాలనే $\lexists[x]$. ఆ పరిమాణీకరణికాన్ని సరిచేసి ప్రకటించాం.}
\item ఇదే తరహా వాదనతో అత్యధిక బ్లాక్ కూడా లేదని తేలుతుంది.
\item బ్లాక్‌ల క్రమం సాంద్రమైనది: $[x]$, $[y]$కన్నా చిన్నదైతే ($x \not\approx y$), వాటి మధ్య రెండింటికీ భిన్నమైన $[z]$ బ్లాక్ ఉంటుంది. $x \prec y$ అనుకుందాం. మునుపటిలాగే $x \oplus
  y$ను రెండు చేత (మిగులు ఉండవచ్చుతో) భాగించవచ్చు; కాబట్టి $u
  \in \Domain{M}$ ఒకటి ఉంటుంది; దానికి $x \oplus y = u \oplus u$ లేదా $x
  \oplus y = u \oplus u \oplus \mathbf{z}^*$. $u$ అనేది $x$, $y$ల సగటు; కనుక వాటి మధ్య ఉంది. $x \oplus y =
  u \oplus u$ అనుకుందాం (ఇతర సందర్భం కూడా ఇదే విధంగా). $u \approx x$ అయితే, ఏదో ప్రమాణ $n$కు:
  \[
  x \oplus y = x \oplus n \oplus x \oplus n,
  \]
  కాబట్టి $y = x \oplus n \oplus n$; అప్పుడు $x \approx y$ అవుతుంది---ఇది పరికల్పనకు విరుద్ధం. కనుక $u \not\approx x$. ఇదే తరహా వాదనతో $u \not\approx y$.
\end{enumerate}
అందువల్ల అప్రమాణ బ్లాక్‌లు కరణీయ సంఖ్యల మాదిరిగా క్రమితమై ఉన్నాయి: అవి చివరి బిందువులు లేని \tetoken{గణనీయ}{!!a{enumerable}} రేఖీయ క్రమాన్ని ఏర్పరుస్తాయి. కనుక సత్య అంకగణితం యొక్క ఏ రెండు \tetoken{గణనీయ}{!!{enumerable}} అప్రమాణ నమూనాలు $\Struct{M}_1$, $\Struct{M_2}$ తీసుకున్నా, $<$, $=$ మాత్రమే ఉన్న భాషకు వాటి సంక్షేపిత నమూనాలు సమరూపాలు. నిజానికి $h$ అనే సమరూపతను ఇలా నిర్వచించవచ్చు: $\Struct{M_1}$, $\Struct{M_2}$ల ప్రమాణ భాగాలు ప్రమాణ నమూనా $\Struct{N}$కు, అందువల్ల ఒకదానికొకటి సమరూపాలు. అప్రమాణ భాగంలోని బ్లాక్‌లు స్వయంగా కరణీయ సంఖ్యల మాదిరిగా క్రమితమవుతాయి; కనుక \olref[dlo]{thm:cantorQ} ప్రకారం అవి సమరూపాలు. ప్రతి బ్లాక్‌లో ఏవైనా మూలకాలను జతచేసి, $\Struct{M_1}$లో $x$ ఉత్తరవర్తి బింబాన్ని $\Struct{M_2}$లో $x$ బింబం ఉత్తరవర్తిగా తీసుకోవడం ద్వారా బ్లాక్‌ల సమరూపతను బ్లాక్‌ల \emph{లోపలికి} విస్తరించవచ్చు. అయితే దీని వల్ల $\mathfrak{M}_1$, $\mathfrak{M}_2$ అంకగణితపు పూర్తి భాషలో సమరూపాలని \emph{తేలదు} (సమరూపత ఎల్లప్పుడూ సంతకం-సాపేక్షమే). ఎందుకంటే $\Domain{M_1}$, $\Domain{M_2}$లపై కూడిక, గుణకారాలను నిర్వచించడానికి అసమరూప మార్గాలున్నాయి. (పూర్తి సిద్ధాంతపు గణనీయ నమూనాల సంఖ్య~2 కాలేదనే వాట్ ప్రసిద్ధ సిద్ధాంతం నుంచీ ఇది వస్తుంది.)

\begin{prob}
అంకగణితపు అప్రమాణ నమూనాలో అత్యధిక బ్లాక్ ఉండదని చూపండి.
\end{prob}

\begin{prob}
$\Lang{L}$ అనేది $<$ను ($\eq$తో పాటు) తన ఏకైక \tetoken{విధేయంగా}{!!{predicate}} కలిగిన మొదటిస్థాయి \tetoken{భాష}{!!{language}} అనుకుందాం; $\Struct{N} = (\Nat,
<)$ అనుకుందాం. $\Struct{N}$ యొక్క ప్రతి పరిమిత లేదా సహపరిమిత ఉపసమితిని నిర్వచించవచ్చు. $\Struct{N}$లో నిర్వచనయోగ్య ఉపసమితులు \emph{ఇవే మాత్రమే} అని చూపండి.

(సూచన: ముందుగా $\Obj{prc}(x,y)$ అనేది ``$x$, $y$కు తక్షణ పూర్వవర్తి'' అని సంక్షిప్తంగా చెప్పే $\Lang{L}$-సూత్రం అనుకుందాం:
\[
x<y \land \lnot \lexists[z][(x<z \land z < y)].
\]
ఇప్పుడు $\Struct{N}$ యొక్క ఏ నిర్వచనయోగ్య ఉపసమితికైనా $\Lang{L}$లో $!A(x)$ అనే సూత్రం ఉంటుంది. అటువంటి ఏ $!A$కైనా ఈ $!D$ వాక్యాన్ని పరిశీలించండి:
\[
\lexists[x][\lforall[y][\lforall[z][((x<y \land x<z \land \Obj{prc}(y,z)
  \land !A(y)) \lif !A(z))]]].
\]
$\Sat{N}{!D}$ కావడానికి, కావలసినదీ చాలునదీ $!A$ నిర్వచించే $\Struct{N}$ ఉపసమితి పరిమితమైనది లేదా సహపరిమితమైనది అని చూపండి.

ఇప్పుడు $\Struct{M}$ అనేది $\Struct{N}$కు ప్రాథమిక తుల్యమైన అప్రమాణ నమూనా అనుకుందాం. $a \in \Domain{M}$ అప్రమాణమైతే, $b, c \in \Domain{M}$లను $a$కన్నా పెద్దవిగా తీసుకుని, $b$ అనేది~$c$కు తక్షణ పూర్వవర్తి అనుకుందాం. అప్పుడు $h(b)=c$ అయ్యే $\Domain{M}$ యొక్క స్వయంసమరూపత $h$ ఒకటి ఉంటుంది (ఎందుకు?). కాబట్టి $b$, $!A$ను సంతృప్తి పరిస్తే $c$ కూడా పరుస్తుంది (ఎందుకు?). దీని వల్ల $!D$, $\Struct{M}$లో సత్యం; అందువల్ల $\Struct{N}$లోనూ సత్యమే. కానీ అప్పుడు $!A$ నిర్వచించే $\Struct{N}$ ఉపసమితి పరిమితమైనది లేదా సహపరిమితమైనది.
\end{prob}

\end{document}
